\chapter{Appendix}
\label{c:appendix}

%%%%%%%%%%%%%%%%%%%%%%%%%%%%%%%%%%%%%%%%%%%%%%%%%%%%%%%%%%%%%%%%%%%%%%%%%%%%%%%%
\section{Supporting equations}
\subsection{Derivation - Calculate fermionic ecitation partner state}\label{subsec:derivation---calculate-fermionic-ecitation-partner-state}
We start from the equation
\begin{equation}
    \ket{\pm} = \frac{1}{\sqrt{2}} \left( \ket{o^p_q} \pm \ket{o^q_p} \right)
\end{equation}~\cite[Appendix - Eigenstates of fermionic generators]{kottmann2020feasible}
``where $\ket{o^p_q}$ denotes all configurations where the $p$ orbitals are occupied and the $q$ orbitals are unoccupied with all combinations allowed for all other orbitals.''~\cite[Chapter II.E.1]{kottmann2020feasible}

Adding and Subtracting these two equations from each other leads us to:
\begin{align}
    \ket{o^p_q} = \frac{1}{\sqrt{2}} \left( \ket{+} + \ket{-} \right), \\
    \ket{o^q_p} = \frac{1}{\sqrt{2}} \left( \ket{+} - \ket{-} \right).
\end{align}

Starting with $\ket{o^q_p}$, we apply the Generator $G$ to the expression:
\begin{equation}
    G\ket{o^q_p} = G \left[ \frac{1}{\sqrt{2}} \left( \ket{+} - \ket{-} \right) \right] = \frac{1}{\sqrt{2}} \left( G\ket{+} - G\ket{-} \right)
\end{equation}
Here, we can apply the eigenvalues of the generator, which are $\lambda_+ = i$ and $\lambda_- = -i$:
\begin{equation}
    G\ket{oqp} = \frac{1}{\sqrt{2}} \left( i\ket{+} - (-i)\ket{-} \right)
\end{equation}
Substituting the before calculated superposition of $\ket{o^p_q}$ we get a equation that already nearly looks like our goal:
\begin{equation}
    G\ket{oqp} = i \left[ \frac{1}{\sqrt{2}} \left( \ket{+} + \ket{-} \right) \right] = i\ket{opq}
\end{equation}
Finally, we multiply both sides with $-i$ and map $\ket{\boldsymbol{i}} = \ket{o^q_p}$, $\ket{\boldsymbol{i'}} = \ket{o^p_q}$.
We introduce $c$ as a real scalar $c = \pm1 = -i \times (\pm i)$ which leads us exactly to:
\begin{equation}
    -iG\ket{\boldsymbol{i}} = c\ket{\boldsymbol{i'}}
\end{equation}

\section{C++ Implementation}
\subsection{Fermionic Excitation}\label{subsec:appendix-fermionic-excitation}

\begin{minted}{cpp}
State apply_fermionic_excitation(const State& state, const FermionTerm& term, double theta) {
    if (state.empty())
        throw std::invalid_argument("A state cannot be empty");

    const std::complex<double> w = term.weight;
    const double abs_w = std::abs(w);
    if (abs_w == 0.0) return state;

    const double cos2 = std::cos(abs_w * theta / 2.0);
    const double sin2 = std::sin(abs_w * theta / 2.0);

    const std::complex<double> i_unit(0.0, 1.0);
    const std::complex<double> factor_a  = -i_unit * w / abs_w;
    const std::complex<double> factor_ad = -i_unit * std::conj(w) / abs_w;

    State new_state;
    new_state.reserve(state.size() * 2);

    for (const auto& [basis_state, coeff] : state) {
        // ---- Nullspace check ----
        bool a_acts  = true;  // A|i⟩ ≠ 0  (all creation empty, all annihilation occupied)
        bool ad_acts = true;  // A†|i⟩ ≠ 0 (all creation occupied, all annihilation empty)

        for (int idx : term.creation_idx) {
            bool bit = (basis_state >> idx) & 1ULL;
            if (bit) a_acts = false;   // creation idx occupied → A can't create
            else ad_acts = false;      // creation idx empty → A† can't annihilate
        }
        for (int idx : term.annihilation_idx) {
            bool bit = (basis_state >> idx) & 1ULL;
            if (bit) ad_acts = false;  // annihilation idx occupied → A† can't create
            else a_acts = false;       // annihilation idx empty → A can't annihilate
        }

        if (!a_acts && !ad_acts) {
            // Nullspace: just copy the basis state unchanged
            new_state[basis_state] += coeff;
            continue;
        }

        // Active space: apply rotation
        int phase = 1;
        uint64_t new_basis_state = basis_state;

        const std::vector<int>& annihilate_idx = a_acts ? term.annihilation_idx : term.creation_idx;
        const std::vector<int>& create_idx     = a_acts ? term.creation_idx : term.annihilation_idx;

        // Apply annihilation
        for (int idx : annihilate_idx) {
            uint64_t mask = (1ULL << idx) - 1;
            if (__builtin_popcountll(new_basis_state & mask) % 2 != 0) {
                phase = -phase;
            }
            new_basis_state ^= (1ULL << idx);
        }

        // Apply creation
        for (int idx : create_idx) {
            uint64_t mask = (1ULL << idx) - 1;
            if (__builtin_popcountll(new_basis_state & mask) % 2 != 0) {
                phase = -phase;
            }
            new_basis_state ^= (1ULL << idx);
        }

        // Update amplitudes
        std::complex<double> factor = a_acts ? factor_a : factor_ad;

        new_state[basis_state]       += coeff * cos2;
        new_state[new_basis_state]   += factor * std::complex<double>(phase, 0.0) * sin2 * coeff;
    }

    return new_state;
}
\end{minted}

\subsection{fSWAP}\label{subsec:appendix-fswap}
\begin{minted}{cpp}
State apply_fswap(const State& state, int i, int j) {
    if (i == j) return state;
    if (state.empty())
        throw std::invalid_argument("A state cannot be empty");

    const uint64_t mask_i = 1ULL << i;
    const uint64_t mask_j = 1ULL << j;

    State new_state;
    new_state.reserve(state.size());

    for (const auto& [basis_state, coeff] : state) {
        bool bit_i = (basis_state >> i) & 1ULL;
        bool bit_j = (basis_state >> j) & 1ULL;

        if (bit_i == bit_j) {
            // |00⟩ → |00⟩, |11⟩ → -|11⟩
            double phase = bit_i ? -1.0 : 1.0;
            new_state[basis_state] += phase * coeff;
        } else {
            // |01⟩ ↔ |10⟩, swap the two bits
            uint64_t swapped = basis_state ^ mask_i ^ mask_j;
            new_state[swapped] += coeff;
        }
    }

    return new_state;
}
\end{minted}

\subsection{Expectation value}\label{subsec:apendix-expectation-value}

\begin{minted}{cpp}
    std::complex<double> expectation_value_fermionic_term(const State &phi, const State &psi, const FermionTerm &term) {
    // Init
    std::complex<double> term_overlap = 0.0;
    //Loop
    for (const auto& [basis_state_psi, coeff_psi] : psi) {
        uint64_t new_basis_state = basis_state_psi;
        int phase = 1;
        bool skip = false;

        // apply Annihilation
        for (int idx : term.annihilation_idx) {
            // Check of the orbital is already empty
            if ((new_basis_state >> idx & 1ULL) == 0) {
                skip = true;
                break;
            }

            uint64_t mask = (1ULL << idx) - 1;
            if (__builtin_popcountll(new_basis_state & mask) % 2 != 0) {
                phase = -phase;
            }
            new_basis_state ^= (1ULL << idx);
        }
        if (skip) continue;

        // apply Creation
        for (int idx : term.creation_idx) {
            // Check of the orbital is already occupied
            if ((new_basis_state >> idx & 1ULL) == 1) {
                skip = true;
                break;
            }
            uint64_t mask = (1ULL << idx) - 1;
            if (__builtin_popcountll(new_basis_state & mask) % 2 != 0) {
                phase = -phase;
            }
            new_basis_state ^= (1ULL << idx);
        }
        if (skip) continue;

        // Compute overlap
        auto it = phi.find(new_basis_state);
        if (it != phi.end()) {
            std::complex<double> coeff_phi = std::conj(it->second);
            term_overlap += coeff_phi * (term.weight * std::complex<double>(phase, 0.0) * coeff_psi);
        }
    }

    return term_overlap;
}

std::complex<double> expectation_value_fermionic(const State &phi, const State &psi, const std::vector<FermionTerm> &operator_sum) {
    std::complex<double> result = 0.0;

    for (const auto &term: operator_sum) {
        result += expectation_value_fermionic_term(phi, psi, term);
    }
    return result;
}
\end{minted}